\begin{example}[A two-stage cone with no common displayed grading]
\label{ex:p061-iterated-weighted-cone-no-common-grading}
Assume \(\operatorname{char}K\ne3\), begin with \(B_0=K\), and put
\[
 B_1=K\langle x,y,z\rangle/(x^3+y^3+z^3),
\]
\[
 B_2=B_1\langle U,V,W\rangle/
       \bigl(UV-(x+y^2)W^2\bigr).
\]
The first equation is homogeneous of degree \(3\) in \(x,y,z\), and the
second is homogeneous of degree \(2\) in the new variables \(U,V,W\), whose
coefficient \(x+y^2\) is allowed to vary on the singular first-stage base.
Thus
\Cref{cor:p061-iterated-relative-weighted-cone-closure} shows that every
classically pluriharmonic germ on \(\mathcal M(B_2)\) at the composite
zero-section point is locally \(\mathbf R\)-PL and
Gubler--Jell--Rabinoff harmonic.

The two displayed equations cannot simultaneously be homogeneous for a
single assignment of strictly positive weights to all six variables.
Indeed, the first equation forces
\(\operatorname{wt}(x)=\operatorname{wt}(y)=\operatorname{wt}(z)\), whereas
the two coefficient terms in the second force
\(\operatorname{wt}(x)=2\operatorname{wt}(y)\).  This only separates the
displayed two-stage presentation from the one-stage weighted presentation;
it is not an assertion that the analytic germ admits no other one-stage
presentation or is intrinsically new.
\end{example}
